\section{Discussion}
\label{sec:discussion}

\para{Forgetting is event-driven.}
The hypothesis pictured forgetting as a smooth process in training time with a task-type-specific rate.
The trajectories show something else: a switch of training distribution causes a fast, bounded shift, after which continued training on the same new task does little further harm.
A ``rate per step'' is then not a natural quantity; ``loss per switch, by pair'' is.
For practice, this means that evaluating shortly after a switch already shows most of the damage, and that training longer on the new task adds little further loss in our setting.

\para{Why does a power law appear in averages?}
The \llama average is well fit by a power law with a small exponent, and so would be many decelerating curves.
The equal-parameter comparison on individual trajectories and the step-resolution experiment both favor saturation.
This matches the known result that averaging can produce power-like curves from components that are not power laws~\cite{anderson1997artifactual}.
We did not run a simulation showing that the observed mixture of steps reproduces the averaged power-law fit, so this remains the most economical reading and not a demonstrated mechanism.

\para{The two models disagree.}
\llama and \qwen disagree on which tasks are fragile (\llama: \tlogic and \tmt; \qwen: \tarith and \tgsm), on whether forgetting accumulates with lag, and on the category ratio.
Any claim of a universal reasoning-vs.-pattern constant would need many more tasks and models than we have.

\para{Normalization matters.}
The same \llama data give a category ratio from 0.88 to 1.47 depending on the measure, and the same \qwen data from 1.2 to 5.9.
Loss-normalized retention breaks down when fine-tuning barely changes a task's loss.
Cross-task forgetting numbers quoted without their normalization are not comparable.

\subsection{Limitations}
\label{sec:limitations}

\para{Scale and method.}
We study 1B and 1.5B models with LoRA at $r{=}16$, one learning rate, 250 steps per task, and no replay or regularization.
We planned a full-fine-tuning ablation and cancelled it before any of its runs started because of GPU contention and time, so nothing here speaks to full fine-tuning or to larger models.

\para{Tasks.}
We use six tasks, three per category, with a debatable category assignment.
Tasks are the real unit of generalization for H2, and there are too few for a powerful test.
Fine-tuning hurt \tgsm accuracy and did not improve \tarith (both models) or \tmt (\qwen).
Forgetting of something that was barely learned is hard to define, and the gain-normalized results rest on 3--4 tasks.

\para{Statistical power.}
Each test set has 200 items, so single accuracy differences below about 7 points are within sampling error; our conclusions rest on aggregates over runs.
The two \llama blocks differ in both seed and order set, so we cannot separate the two.
\qwen has one block and one joint-baseline seed.

\para{Onset experiment.}
It uses one seed, six pairs chosen partly because they showed strong interference in the main runs, \llama only, and losses only.
The S-shaped start coincides with the learning-rate warm-up, and we did not test it without warm-up.

\para{Curve fitting.}
All four pre-registered models are monotone and start at 1, while real trajectories are non-monotone in over 20\% of transitions.
AICc treats the 51 points of a mean curve as independent, which they are not, so the \daicc magnitudes overstate certainty; the win counts are the safer evidence.

\para{Exploratory analyses.}
We added the clipped loss measure, the within-stage fits, the equal-parameter comparison, the sensitivity checks, and the onset experiment after seeing the main results.
The format-validity checks are simple string rules that we validated only by spot checks.

\para{Reproducibility.}
Runs are seeded, but our out-of-memory handling changes micro-batch sizes, so floating-point summation order can differ between repeats.
We do not expect bit-level reproduction and did not repeat a training run to measure the variation.
The analysis scripts are deterministic.

\para{Literature.}
We did not run a systematic literature search for this study.
Our positioning against prior work in \secref{sec:related} reflects our general knowledge of the cited papers, which we did not re-read for this study.
